\chapter{Introduction}
\label{chap:intro}

\section{Motivation}

Modern data pools are rarely purely relational. A movie catalogue such as
IMDb, for instance, pairs a structured relation -- title, release year,
genre, average rating, cast -- with a large body of free-text evidence:
plot summaries and, above all, user reviews. A realistic information need
over such a pool routinely mixes both modalities in a single question:

\begin{quote}
\emph{``Which movies released in 1998 have reviews expressing an overall
negative, critical, or strongly unfavorable opinion of the movie?''}
\end{quote}

The first half of this question -- \texttt{year = 1998} -- is a textbook
relational predicate: it is answered exactly, cheaply, and instantly by a
\texttt{WHERE} clause against an indexed column. The second half --
\emph{``expressing an overall negative, critical, or strongly unfavorable
opinion''} -- is a semantic predicate over unstructured text that no
amount of indexing or exact-match SQL can express. It requires either a
learned text classifier or, at the level of accuracy and flexibility we
target here, a large language model (LLM) invoked as a query operator.

The engineering problem this report addresses is therefore not ``can an
LLM answer this question'' -- it clearly can, given enough context and
enough calls -- but rather: \emph{how do we build a query engine that
answers such questions correctly, while keeping the number of LLM calls
and the wall-clock latency as low as possible?} This is a cost-vs-quality
optimization problem with the same flavor as classical query optimization
(operator ordering, selectivity estimation, pushdown), except that one of
the ``operators'' is a probabilistic, expensive, non-relational function:
the LLM.

\section{Problem Statement}
\label{sec:intro-problem}

Formally, we are given:
\begin{itemize}
  \item a structured relation $R$ with schema
        $(\mathit{id}, \mathit{year}, \mathit{genre}, \mathit{rating},
        \mathit{cast}, \dots)$, $|R| \approx 25{,}000$ rows;
  \item a free-text relation $D(\mathit{id}, \mathit{review})$, in a
        one-to-one (or one-to-few) correspondence with $R$ on
        $\mathit{id}$;
  \item a natural-language query $Q$ that decomposes (possibly only
        implicitly) into a structured predicate $\sigma_S$ over $R$'s
        columns and a semantic predicate $\sigma_U$ over $D$'s text,
        joined conjunctively: $Q \equiv \sigma_S \wedge \sigma_U$.
\end{itemize}
The desired output is the JSON-serializable result set
$\{\, r \in R \bowtie D : \sigma_S(r) \wedge \sigma_U(r) \,\}$, and the
desired \emph{system} property is that it produce this set (i) as close
as possible to a hand-verified gold answer set (measured by precision,
recall, F1), (ii) with as few LLM calls as possible, and (iii) with as
low wall-clock latency as possible. These three objectives are in
tension, and the remainder of this report is an account of how far a
principled combination of classical query-optimization ideas (predicate
pushdown, operator reordering, selectivity-aware short-circuiting) and
modern LLM-serving ideas (batched invocation, confidence calibration,
cascaded inference) can push the achievable trade-off surface.

All code, prompts, and benchmark data referenced in this report are
drawn from the project repository at
\url{https://github.com/AnnaRemi/ENSIMAG-AI-M2-final-project}, branch
\texttt{agent/consolidate-four-method-benchmarks}. Experiments
(Chapter~\ref{chap:experiments}) are run against controlled,
$100$-row-per-question subdatasets drawn from the $\sim$25{,}000-row
canonical corpus above, rather than against the full corpus directly,
so that every question has a known, auditable gold set of a fixed size
(Section~\ref{sec:10q-suite}).

\section{What Existed Before This Project}

Two execution strategies were already implemented and benchmarked before
the work reported here, and both are treated as \emph{baselines} rather
than as strawmen: they are each the right choice for some regime of the
problem, and understanding exactly why is what motivates the cascaded
design of Chapter~\ref{chap:cascading}.

\begin{description}
\item[\textbf{SUQL baseline.}] Following the design of SUQL~\cite{suql},
  the structured predicate $\sigma_S$ is compiled to SQL and evaluated
  first, and the semantic predicate $\sigma_U$ is evaluated by issuing
  one \answerfn{} call per surviving row to an LLM. Quality is high
  because every row is judged individually and in full context, but the
  number of LLM calls equals $|R_{\sigma_S}|$ -- the size of the
  structurally filtered relation -- so the strategy does not scale
  gracefully when the structured predicate is not very selective.
\item[\textbf{Trummer baseline.}] The structured relation and the review
  relation are kept as two \emph{separate} collections -- deliberately,
  by design, with \emph{no} SUQL-style structured pushdown -- and the
  LLM itself is asked to join a block of one collection against a block
  of the other while simultaneously applying the semantic predicate, an
  adaptive block-pair join over the \emph{entire} candidate pool. This
  is far cheaper in calls and wall time than row-wise judgment, but is
  measurably less accurate: the model must recover row identity from
  block text and perform the join itself, and errors (id drift,
  incomplete coverage of a block's decisions) compound without any
  structural narrowing of the pool it has to cover.
\end{description}

Both baselines are described in full, with architecture diagrams and
pseudocode, in Chapter~\ref{chap:architecture}.

\section{Contributions of This Report}

\begin{enumerate}
  \item \textbf{An apples-to-apples empirical comparison} of row-wise
        predicate pushdown and adaptive block-join execution, on ten
        held-out, genre-diverse questions drawn from the canonical
        $\sim$25k-row IMDb corpus, each run ten times per method, with
        matched result sets and a shared metric suite (precision,
        recall, F1, wall time, LLM calls per question) --
        Chapter~\ref{chap:experiments}.
  \item \textbf{A two-level cascaded execution architecture}, applicable
        on top of either baseline, in which a cheap first-stage scorer
        resolves the majority of units of work (rows or batches)
        immediately, and only escalates the units whose Stage~1 output
        falls inside a calibrated ambiguity band to the expensive
        second-stage operator -- Chapter~\ref{chap:cascading}.
  \item \textbf{A closed-form Bayesian calibration rule} for the
        escalation band, derived from a Beta posterior over the cheap
        stage's empirical recall/precision on a held-out calibration
        set, following the backend-agnostic half of Stretto's
        contribution~\cite{stretto} -- Section~\ref{sec:beta-posterior}.
        This rule is fully implemented for SUQL~V1; Trummer~V1
        currently uses a simpler point-estimate threshold search, and
        unifying the two is identified as near-term future work
        (Section~\ref{sec:trummer-v1}).
  \item \textbf{An operator-reordering rule}, inherited from ELEET's
        cascade optimizer~\cite{eleet} and from classical
        selectivity-ordered filter placement, used to decide in which
        order the structural filter, the batched semantic pass, and the
        verification step should run -- Section~\ref{sec:operator-ordering}.
  \item \textbf{A formal proof that calibrated cascading beats an
        all-expensive baseline for large enough candidate pools} -- both
        in expectation (a finite, closed-form crossover point $n^\ast$)
        and, via a Hoeffding concentration bound, with probability
        approaching $1$ exponentially fast in the pool size -- together
        with a worked instantiation from the measured system that
        \emph{predicts} the empirical asymmetry between SUQL~V1 and
        Trummer~V1 rather than merely describing it after the fact --
        Section~\ref{sec:theory-cost}.
  \item \textbf{A full empirical characterization} of the four resulting
        design points (SUQL baseline/V1, Trummer baseline/V1), showing
        that the cascaded batched-join design (Trummer~V1) dominates on
        cost, F1, \emph{and} worst-case robustness across the ten
        held-out questions -- Chapter~\ref{chap:experiments}.
  \item \textbf{An honest account of engineering pitfalls} encountered
        along the way -- an off-by-quantile error in the Bayesian bound,
        hardcoded gold-standard counts, a Trummer baseline that turned
        out (on inspection of the real implementation) to skip
        structural pushdown entirely rather than merely defer it, and a
        missing early-termination optimization for \texttt{LIMIT}
        queries -- because these are as much a part of building a
        reliable LLM-as-operator system as the architecture itself.
\end{enumerate}

\section{Report Roadmap}

Chapter~\ref{chap:background} reviews the three pieces of prior work this
project builds most directly on -- SUQL, ELEET, and Stretto -- and fixes
the formal query model used throughout. Chapter~\ref{chap:architecture}
describes the two baseline architectures in detail, with diagrams,
pseudocode, and the exact prompts used, including a correction of an
earlier draft's mischaracterization of the Trummer baseline's join
structure. Chapter~\ref{chap:cascading} is the technical core of the
report: it develops the two-level cascade, its calibration procedures,
the operator-ordering rule, and -- new in this revision -- a full
theoretical cost analysis proving when and why calibrated cascading
beats an all-expensive baseline, mathematically and in pseudocode, for
both the SUQL~V1 and Trummer~V1 variants. Chapter~\ref{chap:experiments}
presents the ten-question experimental setup and results.
Chapter~\ref{chap:conclusion} concludes and lays out the remaining
roadmap, in particular the KV-cache-aware execution ladder (Stretto's
second contribution) that is deferred pending a backend migration, and
the near-term engineering changes the theory of
Chapter~\ref{chap:cascading} points to directly. Appendix~\ref{chap:appendix}
collects the full system prompts, verbatim from the repository, and the
actual commands used to reproduce every experiment in this report.
